\section{Adelic classes and their integral polynomial coordinates}
\label{acl-section}

The coefficient estimates in Section~\ref{am-section} concern the
actual class lattices, rather than arbitrary lattices with the same
rank. We identify these objects explicitly. The argument below also
gives a short proof of the adelic existence statement needed at the
initial step of the construction. It uses the unramified quadratic
extension and the disjoint selected primes that occur in the
source~\cite{alpo-uniform-source}. Neither strong approximation for
the unitary group nor a class-number formula is required.

\subsection{Unit norms and integral orthonormal bases}

\begin{lemma}[Norms in an unramified quadratic extension]
\label{acl-local-norm}
Let $A$ be a complete discrete valuation ring with finite residue
field $\kappa$, and let $B$ be the ring of integers in an unramified
quadratic extension of its fraction field. Write $\bar x$ for the
nonidentity automorphism. Then the residue trace
$\kappa_B\to\kappa$ is surjective, and
\[
                         N_{B/A}:B^\times\longrightarrow A^\times
\]
is surjective. These statements include residue characteristic two.
\end{lemma}
\begin{proof}
Write $q=|\kappa|$. The trace is $x\mapsto x+x^q$ on
$\mathbb F_{q^2}$. In odd characteristic its value at $1$ is
nonzero. In characteristic two its value at any
$x\notin\mathbb F_q$ is nonzero, since $x^q\ne x$.
It is therefore a nonzero linear map from a two-dimensional
$\mathbb F_q$-space to a one-dimensional space, hence is onto.

The multiplicative group of a finite field is cyclic. One elementary
proof is to let $e$ be the exponent of this finite abelian group.
Combining elements of maximal prime-power orders gives an element
of order $e$. Every group element is a root of $X^e-1$, so the
group has at most $e$ elements; its exponent also divides its
order. The two numbers are equal. It follows that
$x\mapsto x^{q+1}$ maps $\mathbb F_{q^2}^\times$ onto
$\mathbb F_q^\times$.

Choose a uniformizer $\varpi$ of $A$, which is also a uniformizer
of $B$. Given $a\in A^\times$, first choose $u_1\in B^\times$
with $N(u_1)\equiv a\pmod\varpi$. Suppose inductively that
$N(u_j)\equiv a\pmod{\varpi^j}$. For $t\in B$,
\[
 N(1+\varpi^jt)
  =1+\varpi^j(t+\bar t)+\varpi^{2j}t\bar t
  \equiv1+\varpi^j\operatorname{Tr}(t)
                                     \pmod{\varpi^{j+1}}.
\]
Surjectivity of the residue trace permits a choice of $t_j$
such that $u_{j+1}=u_j(1+\varpi^jt_j)$ has
$N(u_{j+1})\equiv a\pmod{\varpi^{j+1}}$.
The $u_j$ converge in $B^\times$, because their successive
ratios belong to $1+\varpi^jB$. Their limit has norm $a$.
No division by two occurs.
\end{proof}

\begin{proposition}[Integral orthonormal bases]
\label{acl-local-basis}
Let $A$ be as above, and let $B$ be either its unramified quadratic
extension or $A\times A$, with the standard involution. Every
unimodular Hermitian form on a free $B$-module of rank $n$ has a
basis $e_1,\ldots,e_n$ with $h(e_i,e_j)=\delta_{ij}$.
\end{proposition}
\begin{proof}
First suppose $B$ is a discrete valuation ring. The Gram matrix
in an arbitrary basis belongs to $\operatorname{GL}_n(B)$.
If some diagonal entry is a unit, the corresponding vector has
unit norm. Otherwise some off-diagonal entry $a=h(e_i,e_j)$
is a unit. The reduction of
\[
 h(e_i+te_j,e_i+te_j)
   =h(e_i,e_i)+\operatorname{Tr}(ta)+t\bar t\,h(e_j,e_j)
\]
is nonzero for a suitable $t\in B$, by
Lemma~\ref{acl-local-norm}. Thus a vector of unit norm exists
also in this case. Multiplying it by a unit whose norm is the
inverse of its norm gives a vector $v$ with $h(v,v)=1$.

The integral projection $x\mapsto v h(v,x)$ splits off $Bv$.
Its orthogonal complement is a direct summand, hence free over
the discrete valuation ring. The Gram determinant shows that
the restricted form is again unimodular. Induction on $n$
proves the assertion, including the empty basis when $n=0$.

If $B=A\times A$, a free $B$-module is the product of two
free $A$-modules of rank $n$. A Hermitian form amounts to a
bilinear pairing between them, with the reverse pairing
determined by symmetry. Unimodularity says that this pairing
is perfect. Choose a basis on one factor and its dual basis
on the other. The corresponding $B$-basis has Gram matrix $1$.
\end{proof}

\subsection{Patching lattices and prescribing split-prime data}

Let $E/F$ be a CM quadratic extension, unramified at every finite
prime. Put $\mathcal O=\mathcal O_E$, and fix
\[
 V=E^n,\qquad h(x,y)=\sum_{i=1}^n\bar x_i y_i,
 \qquad N_0=\mathcal O^n,\qquad n\geq1.
\]
Let $U$ be its unitary group over $F$. At a finite place $v$ of
$F$, write $E_v=E\otimes_FF_v$ and $\mathcal O_v$ for its
maximal compact subring. The group $U(\mathbb A_{F,f})$ consists
of tuples in $U(F_v)$ preserving $(N_0)_v$ at all but finitely
many places.

\begin{lemma}[Patching full lattices]
\label{acl-patching}
Suppose a full lattice $M_P\subset E_P^n$ is prescribed at
each prime $P$ of $\mathcal O$, with $M_P=\mathcal O_P^n$
for all but finitely many $P$. Here $\mathcal O_P$ denotes the
completed discrete valuation ring. There is a unique full
$\mathcal O$-lattice $M\subset E^n$ with these completions:
\[
                         M=E^n\cap\prod_PM_P.
\]
In particular, for $g\in U(\mathbb A_{F,f})$,
\[
                         N_g=V\cap g\widehat N_0
\]
is unimodular and has completion $g_v(N_0)_v$ at every $v$.
\end{lemma}
\begin{proof}
Choose a nonzero rational integer $a$ such that, at every prime,
\[
                 a\mathcal O_P^n\subseteq M_P
                              \subseteq a^{-1}\mathcal O_P^n.
\]
Only finitely many local conditions need be imposed when choosing
$a$. The finite $\mathcal O$-module
$Q=a^{-1}\mathcal O^n/a\mathcal O^n$ is annihilated by $a^2$.
The Chinese remainder theorem decomposes it as the direct sum
of its localizations at the primes dividing $a$; these finite
modules are unchanged by completion. Each $M_P$ determines a
submodule of the corresponding local summand of $Q$. Take their
direct sum and its inverse image in $a^{-1}\mathcal O^n$.
The resulting module has exactly the required completions and
contains $a\mathcal O^n$, so it is full and finitely generated.

The same finite quotient shows that a lattice is the intersection
of its completions with $E^n$, proving uniqueness. Hermitian
duality commutes with completion over $F$: in local bases it is
given by the inverse transpose of the conjugate Gram matrix.
Each $g_v(N_0)_v$ is self-dual because $g_v$ is unitary.
Uniqueness therefore gives $N_g=N_g^\#$.
\end{proof}

Fix disjoint finite sets $S,G$ of primes of $\mathcal O$ such that
\begin{equation}\label{acl-primes}
               P\in S\cup G\quad\Longrightarrow\quad
                                      \bar P\notin S\cup G.
\end{equation}
Every selected prime lies over a split prime of $F$: in a quadratic
Galois extension an inert or ramified prime is fixed by conjugation.
Condition~\eqref{acl-primes} also says that no two selected primes
lie over the same prime of $F$.
At each $P\in S$, fix $c_P\in(\mathcal O/P)^\times$ and the
functional $\ell_{0,P}(x)=c_Px_n$ on $N_0$.
At each $Q\in G$, use the standard coordinate frame
$b_{0,Q}:N_0/QN_0\simeq(\mathcal O/Q)^n$.

A decorated lattice consists of a unimodular $N\subset V$, a
nonzero $\mathcal O$-linear map $\ell_P:N\to\mathcal O/P$ for each $P\in S$,
and an $\mathcal O/Q$-linear isomorphism
$b_Q:N/QN\to(\mathcal O/Q)^n$ for each
$Q\in G$. These are exactly the initial-step lattice and residue
data. The associated integral lattice is
\begin{equation}\label{acl-kernel}
                      L_d=\{x\in N:\ell_P(x)=0\text{ for all }P\in S\}.
\end{equation}
The maps on localizations, and their zero extensions when needed,
are uniquely determined by these restrictions, as in
Section~\ref{al-section}.

For $g\in U(\mathbb A_{F,f})$, transport the fixed functionals
and frames by $g^{-1}$ to obtain a decoration $d_g$ of $N_g$.
Local reduction agrees with $N_g/PN_g$ by the patching lemma.
In particular these functionals are nonzero and the frames are
isomorphisms.

\begin{theorem}[Every decorated lattice has an adelic representative]
\label{acl-points}
Every decorated lattice just defined equals $d_g$ for some
$g\in U(\mathbb A_{F,f})$.
\end{theorem}
\begin{proof}
Complete $N$ at every finite place $v$ of $F$. By
Proposition~\ref{acl-local-basis}, there is an isometry
$x_v:(N_0)_v\to N_v$. Outside finitely many places $N_v=(N_0)_v$,
and we take $x_v=1$ there. Hence $(x_v)$ is an adelic unitary
element.

At a selected split prime $v$, the local unitary group is
$\operatorname{GL}_n(F_v)$: after identifying
$E_v=F_v\times F_v$, its two components have the form
$A$ and $(A^{-1})^t$. The stabilizer of $(N_0)_v$ is
$\operatorname{GL}_n(\mathcal O_{F_v})$, whose reduction maps
onto $\operatorname{GL}_n(\kappa_v)$. To see the latter
surjectivity, lift the entries of any invertible residue matrix;
its determinant is a unit, so the lifted matrix is invertible.

If $P\in S$ lies over $v$, the pulled-back map $\ell_Px_v$
is a nonzero row vector over $\kappa_v$. The general linear
group acts transitively on nonzero row vectors: complete a
nonzero linear functional to a basis of the dual space.
Choose an integral local isometry $a_v$ such that
$\ell_Px_va_v=\ell_{0,P}$ modulo $P$.
If $Q\in G$ lies over $v$, choose instead $a_v$ with
$b_Qx_va_v=b_{0,Q}$ modulo $Q$, by lifting the inverse of
the residue matrix of $b_Qx_v$.
There is only one condition at each such $v$, by
\eqref{acl-primes}. Put $g_v=x_va_v$ at these places and
$g_v=x_v$ elsewhere. It has the same completed lattice as $x$,
and the prescribed functionals and frames. Thus $d_g=d$.
\end{proof}

Multiplying each $\ell_P$ by a fixed nonzero residue scalar
changes neither $L_d$ nor any stabilizer. This accounts for the
source's fixed diagonal rescaling of the level-prime decorations.
In particular, if $\pi\bar\pi=T$ and $(\pi,\bar\pi)=1$, then
$\bar\pi$ is a unit at each prime dividing $\pi$, so it is an
allowed choice of the scalar $c_P$. The proof does not assume
that $N$ has a global integral basis.

\subsection{Double cosets, stabilizers, and the inverse level character}

Let $K_0$ be the stabilizer of $\widehat N_0$, and let
$K_S\subseteq K_0$ preserve every $\ell_{0,P}$.
Reduction at the primes in $G$ gives a surjective homomorphism
\[
       r:K_S\longrightarrow\mathcal G,
       \qquad \mathcal G=\prod_{Q\in G}
                                \operatorname{GL}_n(\mathcal O/Q).
\]
Surjectivity follows by the same independent split-prime lifts
used in Theorem~\ref{acl-points}. Fix a subgroup $H\subseteq\mathcal G$
and put $K_H=r^{-1}(H)$ and $K_1=\ker r$.
These groups are compact and open: they impose finitely many
residue conditions inside the compact open group $K_0$.

The group $U(F)\times H$ acts on decorated lattices by
transporting the lattice and functionals with $\gamma\in U(F)$
and postcomposing the residue frames with $h\in H$.
Directly from the definition by inverse pullback,
\begin{equation}\label{acl-action}
 d_{\gamma g}=(\gamma,1)d_g,
 \qquad d_{gk}=(1,r(k)^{-1})d_g\quad(k\in K_S).
\end{equation}

\begin{proposition}\label{acl-double-cosets}
The maps above induce bijections
\[
 U(\mathbb A_{F,f})/K_1\simeq\{\text{decorated lattices}\},
 \qquad
 U(F)\backslash U(\mathbb A_{F,f})/K_H
       \simeq (U(F)\times H)\backslash\{\text{decorated lattices}\}.
\]
These double-coset sets are finite. For a representative $g$,
the stabilizer of $d_g$ in $U(F)\times H$ is isomorphic to
\begin{equation}\label{acl-rational-stabilizer}
                         \Gamma_g=U(F)\cap gK_Hg^{-1}
\end{equation}
by first projection. Its inverse sends $\gamma$ to
$(\gamma,r(g^{-1}\gamma g))$.
\end{proposition}
\begin{proof}
If $d_g=d_{g'}$, their completed lattices agree, so
$k=g^{-1}g'$ preserves $\widehat N_0$. Equality of the
functionals gives $k\in K_S$, and equality of the frames
gives $r(k)=1$. Conversely every such $k$ fixes $d_g$.
Theorem~\ref{acl-points} now proves the first bijection.
Equation~\eqref{acl-action}, together with the surjectivity
of $r$, proves the second.

Finiteness follows from Corollary~\ref{al-finite-decorated}.
For the stabilizer assertion, if
$\gamma\in\Gamma_g$ and $k=g^{-1}\gamma g$, then
\[
                  (\gamma,1)d_g=d_{\gamma g}=d_{gk}
                                      =(1,r(k)^{-1})d_g.
\]
Thus $(\gamma,r(k))$ fixes $d_g$. The converse follows by
the same computation and the first bijection. Once $\gamma$
is specified, the frame isomorphisms force the component in
$H$ uniquely, proving injectivity. These stabilizers are finite
by Lemma~\ref{al-finite-isometries}; at the two selected type
primes they are trivial under the unipotence hypotheses of
Proposition~\ref{al-two-primes}.
\end{proof}

Let $K$ be a number field containing the coefficient embeddings
of $E$ and the values of a character $\chi:H\to K^\times$.
Since $H$ is finite, these values are roots of unity. The
corresponding right level character is
\begin{equation}\label{acl-character}
                            \omega(k)=\chi(r(k))^{-1}.
\end{equation}
This inverse is necessary, as \eqref{acl-action} shows.

\subsection{Additional local bases and the row classes}

The row objects also occur in a description with actual local
bases at the auxiliary character primes. Forgetting a local basis
and remembering its residue frame does not change either the
class set or the stabilizer order. We give the argument explicitly,
since division by the group of basis changes would be incorrect:
that group is generally infinite.

\begin{lemma}\label{acl-free-fibers}
Let $p:\widetilde D\to D$ be a surjective map, equivariant for
a surjective homomorphism $\pi:\widetilde\Gamma\to\Gamma$.
Suppose $\ker\pi$ acts freely and transitively on each fiber
of $p$. Then $p$ induces a bijection of orbit sets, and for
every $x\in\widetilde D$ the homomorphism $\pi$ restricts
to an isomorphism
\[
           \operatorname{Stab}_{\widetilde\Gamma}(x)
                    \simeq\operatorname{Stab}_\Gamma(p(x)).
\]
\end{lemma}
\begin{proof}
If $p(y)=\gamma p(x)$, lift $\gamma$ to $g\in\widetilde\Gamma$.
Then $y$ and $gx$ lie in the same fiber, so $y=kgx$ for
some $k\in\ker\pi$. This proves injectivity on orbit sets;
surjectivity follows from that of $p$. For a stabilizer
element $\gamma$ downstairs, choose a lift $g$ and then $k$
with $kgx=x$. Thus the stabilizer map is onto. Its kernel
is trivial because $\ker\pi$ acts freely on the fiber.
\end{proof}

Let $Q$ be a finite set of additional primes disjoint from
$S\cup G$, and assume that $S\cup G\cup Q$ still contains
no conjugate pair. At $P\in Q$ put
$A_P=\mathcal O_{E,(P)}$, the localization inside $E$.
An enhanced decoration records a basis matrix
$B_P\in\operatorname{GL}_n(E)$ with
$B_PA_P^n=N\otimes_{\mathcal O}A_P$, rather than a residue
frame at $P$. Forget it by remembering
\[
                            b_P(x)=B_P^{-1}x\pmod P.
\]
Every residue frame has such a lift: first choose any basis
of the free $A_P$-module, then lift the required invertible
residue change of basis to $\operatorname{GL}_n(A_P)$.

For each $P\in Q$, choose a subgroup
$H_P\subseteq\operatorname{GL}_n(\mathcal O/P)$ and let
$\widetilde H_P$ be its inverse image in
$\operatorname{GL}_n(A_P)$. The acting groups are
\[
 \widetilde\Gamma=U(F)\times H\times\prod_{P\in Q}\widetilde H_P,
 \qquad
 \Gamma=U(F)\times H\times\prod_{P\in Q}H_P.
\]
The local action is $B_P\mapsto\gamma B_Pk_P^{-1}$.
Its residue frame is transformed by
$b_P\mapsto\bar k_Pb_P\gamma^{-1}$, so the forgetful map
is equivariant for the evident reduction homomorphism.
The kernel acts freely on each fiber: a basis matrix has no
nonidentity right stabilizer. It acts transitively as well,
since two bases with the same residue frame differ by a
unique matrix congruent to $1$ modulo $P$.
Lemma~\ref{acl-free-fibers} therefore applies.

\begin{corollary}[The row class set and its stabilizers]
\label{acl-row-classes}
Take $S=\varnothing$, so $L_d=N=N^\#$. The row descriptions
by enhanced decorations and by residue-decorated unimodular
lattices have canonically bijective class sets and isomorphic
stabilizers. The latter classes are the double cosets for
the compact open subgroup with the prescribed residue groups
at $G\cup Q$. All these sets are finite. The statement also
holds in rank zero: every group and every frame is then
trivial, and there is one zero lattice.
\end{corollary}
\begin{proof}
For positive rank, combine Lemma~\ref{acl-free-fibers} with
Proposition~\ref{acl-double-cosets}, using $G\cup Q$ in place
of $G$. The rank-zero assertion follows directly from the
definitions.
\end{proof}

Characters which factor through the residue groups pass through
this isomorphism with their exact values. If the row character
is denoted by $\chi_{\rm row}$ and the character on the local
basis description by $\chi_{\rm pt}$, the convention in the
source is
\begin{equation}\label{acl-row-character}
              \chi_{\rm pt}(g)
                        =\chi_{\rm row}(\pi g)^{-1}.
\end{equation}
Indeed, the former is the product of the chosen local
characters, while the latter is the product of their inverses.
This is compatible with the contragredient coefficient action
on one side of the period pairing and the covariant action
on the other. The equality of stabilizer orders above shows
that both descriptions have exactly the same weights in
their finite class sums. No index of the principal congruence
groups is introduced.

\subsection{The polynomial representation and the class lattice}

Fix embeddings $\tau:E\hookrightarrow K$ indexed by a finite
set $\Psi$, and dominant integral weights
$\lambda_\tau=(\lambda_{\tau,1},\ldots,\lambda_{\tau,n})$.
Write $a_\tau=\lambda_{\tau,n}$ and
$\nu_{\tau,i}=\lambda_{\tau,i}-a_\tau$. Let $W_\lambda(K)$
be the tensor product of the polynomial spaces already defined
in Section~\ref{coef:section}, with the determinant twists
recorded separately. In a block of variables $X_\tau$, the
covariant representation is
\[
 \rho_\lambda(A)F(X)
       =\left(\prod_\tau\det(A_\tau)^{a_\tau}\right)
                                         F((A_\tau^tX_\tau)_\tau).
\]
The representation used for evaluations on frames is its
contragredient realization
\begin{equation}\label{acl-contragredient}
 R_\lambda(\gamma)F(X)
    =\left(\prod_\tau\det(\tau\gamma)^{-a_\tau}\right)
                                  F(((\tau\gamma)^{-1}X_\tau)_\tau)
    =\rho_\lambda(((\tau\gamma)^{-1})^t)_\tau F.
\end{equation}
For an embedding $\sigma:K\hookrightarrow\mathbb C$, evaluate
$F$ on a tuple of invertible complex frames $Z=(Z_\tau)$ by
\begin{equation}\label{acl-evaluation}
               f_F(Z)=\prod_\tau\det(Z_\tau)^{a_\tau}
                                      \sigma(F)(Z).
\end{equation}
Negative $a_\tau$ are allowed, since the frames are invertible.
The defining identity is
\begin{equation}\label{acl-eval-action}
               f_{R_\lambda(\gamma)F}(Z)
                    =f_F(((\sigma\tau\gamma)^{-1}Z_\tau)_\tau).
\end{equation}

Evaluation is injective. Indeed, its determinant factor never
vanishes. If a polynomial vanishes on all invertible matrices,
fix any matrix $A$ and consider $A+tI$. Its determinant is a
nonzero polynomial in $t$, so the given polynomial vanishes
at $A+tI$ for all but finitely many $t$ and hence also at $t=0$.
Apply this argument to each block. A polynomial over an infinite
field which vanishes on every tuple has zero coefficients, by
induction on the number of variables.

For a full lattice $M\subset E^n$, denote by
$\Delta_{\lambda,\tau_\bullet}^{\mathrm{ctr}}(M)$ the
integral Weyl lattice transported by $R_\lambda$ from a local
basis of $M$ at every coefficient prime. Its independence of
these bases, finite generation, and compatibility with
localization were proved in Section~\ref{coef:section}.
Equivalently, at a prime $w$ of $K$, if $B_\tau$ is a basis
matrix of $M$ at $P_\tau=\tau^{-1}(w)$, this lattice is
\[
 \rho_\lambda\bigl(((\tau B_\tau)^{-1})^t\bigr)_\tau
                             \Delta_{\lambda}^{\mathrm{std}}(\mathcal O_{K,w}).
\]
This is the coefficient lattice in the actual class construction,
with $M=L_d^\#$.

\begin{lemma}[Conjugation in the coefficient lattice]
\label{acl-conjugate-lattice}
For the standard Hermitian form,
\begin{equation}\label{acl-conjugate-lattice-eq}
 \Delta_{\lambda,\tau_\bullet}^{\mathrm{ctr}}(L^\#)
                 =\Delta_{\lambda,\tau_\bullet\circ\bar{\ }}(L).
\end{equation}
The right-hand side is the covariant Weyl lattice evaluated
through the conjugate embeddings. The same identity holds
for the integral dual Weyl lattices.
\end{lemma}
\begin{proof}
Fix a coefficient prime $w$ and put $P_\tau=\tau^{-1}(w)$.
Choose a basis matrix $A_\tau$ of $L$ at the conjugate prime
$\bar P_\tau$. A basis of $L^\#$ at $P_\tau$ is
\[
                             B_\tau=(\bar A_\tau)^{-t}.
\]
This follows by writing the pairing in coordinates:
$h(Ax,y)=\bar x^t\bar A^ty$.
Substitution in the preceding local description gives
\[
        ((\tau B_\tau)^{-1})^t
                     =\tau(\bar A_\tau)
                     =(\tau\circ\bar{\ })(A_\tau).
\]
This is precisely the local description on the right of
\eqref{acl-conjugate-lattice-eq}, including every determinant
twist. Equality at all coefficient primes proves the global
identity. The argument applies to both integral coefficient
lattices because both have the same change-of-basis rule.
\end{proof}

Define $\mathcal F_K$ to be the space of functions $F$ from
decorated lattices to $W_\lambda(K)$ satisfying
\begin{equation}\label{acl-equivariant}
             F((\gamma,h)d)=\chi(h)R_\lambda(\gamma)F(d).
\end{equation}
Its integral lattice is
\[
 \mathcal F_{\mathcal O_K}
   =\{F\in\mathcal F_K:
             F(d)\in\Delta_{\lambda,\tau_\bullet}^{\mathrm{ctr}}(L_d^\#)
                                          \text{ for every }d\}.
\]
The condition is preserved by the action. Indeed,
$L_{(\gamma,h)d}=\gamma L_d$, the change-of-basis rule gives
the required coefficient-lattice transport, and roots of unity
are integral units.

\begin{proposition}[Evaluation at class representatives]
\label{acl-class-evaluation}
Choose representatives $d_c$ of the finite decorated class set,
and let $\Gamma_c$ be their stabilizers. Evaluation gives an
isomorphism
\begin{equation}\label{acl-evaluation-product}
 \mathcal F_K\simeq\bigoplus_c W_\lambda(K)^{\Gamma_c},
 \qquad
 \mathcal F_{\mathcal O_K}\simeq
   \bigoplus_c\left(
        \Delta_{\lambda,\tau_\bullet}^{\mathrm{ctr}}(L_{d_c}^\#)
                           \cap W_\lambda(K)^{\Gamma_c}\right),
\end{equation}
where $\Gamma_c$ acts by $\chi(h)R_\lambda(\gamma)$.
If the stabilizers are trivial, the invariant restrictions disappear.
\end{proposition}
\begin{proof}
An equivariant function takes a $\Gamma_c$-invariant value at
$d_c$, and all its other values are then determined. Conversely,
given invariant vectors $v_c$, define
\[
                     F((\gamma,h)d_c)
                               =\chi(h)R_\lambda(\gamma)v_c.
\]
Two choices differ by an element of $\Gamma_c$, so this is
well-defined. The representation law proves equivariance.
The lattice assertion follows from the transport property just
proved. In particular this argument neither averages over a
stabilizer nor introduces a denominator equal to its order.
\end{proof}

For clarity, localization introduces no such denominator either.
If $M_c$ is any finite collection of full coefficient lattices and
$W_0\subseteq\bigoplus_c W_\lambda(K)$ is a $K$-subspace, then
\begin{equation}\label{acl-localization}
 \left(W_0\cap\bigoplus_cM_c\right)_w
                     =W_0\cap\bigoplus_c(M_c)_w.
\end{equation}
For a vector on the right, each component has an integral
multiple with multiplier outside $w$. The product of the
finitely many multipliers works for every component and remains
outside $w$; it also preserves $W_0$. This proves the reverse
inclusion, while the forward inclusion is immediate. A full
lattice is the intersection of all its localizations: if
$x\notin M$, the proper annihilator of $x+M$ is contained in
some maximal ideal $w$, and $x\notin M_w$.
Consequently the displayed local coefficient conditions are
equivalent to the global lattice in
\eqref{acl-evaluation-product}.

\subsection{The adelic functions and the metric}

Over $\mathbb C$, associate to $F$ the function
\[
                              A_F(g,Z)=f_{F(d_g)}(Z).
\]
Equations~\eqref{acl-action}, \eqref{acl-eval-action}, and
\eqref{acl-equivariant} give exactly
\begin{equation}\label{acl-adelic-law}
 A_F(\gamma gk,((\sigma\tau(\gamma))Z_\tau)_\tau)
                           =\sigma(\omega(k))A_F(g,Z)
 \qquad(\gamma\in U(F),\ k\in K_H).
\end{equation}
Conversely, a function satisfying this law and having the prescribed
polynomial form at each $g$ determines $F$ uniquely.
Indeed, use Theorem~\ref{acl-points} to choose $g$ for each $d$;
another choice differs by $K_1$, whose character is one.
Injectivity of frame evaluation makes the polynomial independent
of the choice. To recover \eqref{acl-equivariant}, lift
$h^{-1}\in H$ to $k\in K_H$ and use
$d_{\gamma gk}=(\gamma,h)d_g$.
The rational subspace consists exactly of those functions whose
class polynomials have coefficients in $\sigma(K)$ and belong
to $\sigma(W_\lambda(K))$. The integral subspace is the one
specified in \eqref{acl-evaluation-product} and
\eqref{acl-localization}.

Give $W_\lambda(\mathbb C)$ its normalized Fock inner product,
as in Section~\ref{coef:section}. A unitary matrix acts
unitarily on the polynomial representation: the standard Gaussian
integral is invariant under its change of variables, and the
determinant character has absolute value one.
Every $\sigma\tau(\gamma)$ is unitary, since complex conjugation
on the CM field induces the given involution at every embedding.
Also $|\sigma\chi(h)|=1$. Hence the class norm is
\begin{equation}\label{acl-class-norm}
                  \|F\|_\sigma^2
                     =\sum_c\frac{\|\sigma(F(d_c))\|_{\mathrm{Fock}}^2}
                                             {|\Gamma_c|},
\end{equation}
independently of the representatives.

This is also the adelic quotient norm when the Haar measure is
normalized by $\operatorname{vol}(K_H)=1$. Indeed, the
double coset represented by $g_c$ is the quotient of $K_H$
by the finite group $g_c^{-1}\Gamma_{g_c}g_c$, and its quotient
mass is $1/|\Gamma_{g_c}|$. The coefficient norm is constant
on this double coset. The stabilizer identification in
Proposition~\ref{acl-double-cosets} gives exactly
\eqref{acl-class-norm}. One may equivalently take this finite
sum as the definition of the algebraic automorphic norm.
For an arbitrary fixed Haar measure $\mu$, the corresponding
integral is $\mu(K_H)$ times \eqref{acl-class-norm}.
Thus the algebraic class norm uses the measure divided by
$\mu(K_H)$; it does not assert that a separately chosen mass
or formal-degree normalization already has $\mu(K_H)=1$.
There is no additional factor $|H|$ in the finite class norm.
To justify the discrete quotient in this interpretation, a compact
open neighborhood meets $U(F)$ in a finite lattice-isometry group.
A smaller neighborhood excludes its nonidentity elements, proving
discreteness. The finite double-coset decomposition then reduces
the quotient integral to the displayed finite sum.

Under the actual two-prime unipotence hypotheses, every
$|\Gamma_c|$ equals one. Combining
\eqref{acl-conjugate-lattice-eq} with
\eqref{acl-evaluation-product} and \eqref{acl-class-norm}
therefore identifies the initial column lattice isometrically with
\begin{equation}\label{acl-actual-class-lattice}
                  \bigoplus_c
                         \Delta_{\lambda,\tau_\bullet\circ\bar{\ }}
                                      (L_{d_c}),
\end{equation}
with the standard coefficient metric on each summand.
The class lattices $L_{d_c}$ are integral kernels in unimodular
Hermitian lattices. Thus the degree and short-dual conclusions
of Section~\ref{al-section}, and the coefficient maps and slope
bound of Section~\ref{am-section}, apply to these actual objects.

\subsection{Necessary hypotheses and exact scope}

Several small examples guard the identifications above.

\begin{itemize}
\item The unramified hypothesis in Lemma~\ref{acl-local-norm}
cannot be dropped. In $\mathbb Q_2(i)/\mathbb Q_2$, a unit norm
is $a^2+b^2$ with exactly one of $a,b\in\mathbb Z_2$ odd,
so it is $1$ modulo $4$. In particular $-1$ is not a unit
norm. The unimodular rank-one forms with Gram entries $1$
and $-1$ are not integrally isometric. This does not affect
the unramified extension used above.
\item Arbitrary simultaneous decorations at conjugate split
primes cannot be prescribed. In rank one over
$\mathbb F_5\times\mathbb F_5$, a unitary matrix is
$(t,t^{-1})$. The two pulled-back functional coefficients
are $t^{-1}$ and $t$, with product one. The pair with
coefficients $1$ and $2$ cannot arise. Condition
\eqref{acl-primes} excludes precisely this obstruction.
\item Conjugation in \eqref{acl-conjugate-lattice-eq} matters.
In $E=\mathbb Q(i)$, let $P=(2+i)$ and take rank one and
weight $1$. Then $P^\#=\bar P^{-1}$, so its contragredient
coefficient lattice is $\bar P$, not $P$. The two ideals
are distinct.
\item Stabilizer invariants cannot be discarded without
triviality. A one-dimensional coefficient space with the
sign action of a stabilizer $\{\pm1\}$ has zero invariant
space, although the unrestricted coefficient space is
one-dimensional. The two-prime argument is used before
asserting the full direct sum in
\eqref{acl-actual-class-lattice}.
\end{itemize}

In the actual construction, the Hermitian spaces have standard
orthonormal frames and the standard integral lattice by their
definition. The class lattice is the inverse-coordinate image
of the selected decorated representative's kernel. Passing back
to coordinates gives exactly that kernel, and the frame
determinant is the ordinary matrix determinant. The selected
primes satisfy \eqref{acl-primes}: primes over $T$ divide
$\pi_T$ and do not divide $\bar\pi_T$, the type primes use
the prescribed conjugate half of the divisors of $g_1g_2$,
and the distinct rational primes keep these sets disjoint.
These are supplied hypotheses of the prime construction,
not new existence assumptions.

At the initial step the auxiliary coupling set is empty, the
index set has one element, and the reference character has no
auxiliary character primes. Thus no additional class label or
local-basis choice survives in \eqref{acl-actual-class-lattice}.
For the row step the coupling set is again empty; its extra
local bases are handled by Corollary~\ref{acl-row-classes}.

The present section proves the initial-step and row class
comparisons, the initial coefficient descriptions, and its
metric, together with the exact equality of stabilizer weights
needed for the row pairing. More elaborate intermediate steps also carry auxiliary
cosets and local character data. Their comparison with the
intermediate adelic level charts requires the corresponding
additional source arguments. The period-map valuation bound
and the global theta identities are not consequences of this
class comparison.


